\section{Introduction}\label{sec:intro}

The physics described by non-abelian gauge theories
can be studied in many ways. Depending on the
context and the questions to be answered,
the desired information is extracted from
correlation functions of local fields,
expectation values of Wilson loops or the Schr\"odinger
functional, for example.

As explained in refs.~\cite{WilsonFlow,Villasimius},
the gradient flow provides further opportunities
to probe these theories
and allows some of their otherwise elusive properties
to be understood
(in lattice gauge theory, the gradient flow is referred to
as the Wilson flow,
while in the mathematical literature it
is commonly known as the Yang--Mills gradient flow).
Evidently, physically meaningful probes must be safe of ultra-violet
divergences or must be such that these can be canceled by
a well-defined renormalization procedure.

In the case of the gradient flow,
there is some evidence that the gauge field generated by the
flow does not require renormalization \cite{WilsonFlow,Villasimius},
but a formal proof of the absence of ultra-violet divergences at
positive flow time has not been given so far.
The aim of the present paper is to fill this gap
through an all-order analysis of the flow in perturbation theory.
With respect to the closely related case of the renormalization of
the Langevin equation discussed by
Zinn--Justin and Zwanziger \cite{Zinn,ZinnZwanziger},
there are two important differences, one being the absence of
the noise term in the flow equation and the other the fact that
the initial distribution of the gauge field
(which is given by the functional integral of the theory considered)
is not ignored.

The perturbative analysis presented in this paper applies to
renormalizable gauge theories with any compact simple gauge group
and any matter multiplet.
However, in order to simplify the discussion as much as
possible, only the pure $\SUn$ gauge theory with dimensional
regularization will be considered, the generalization
to other cases being straightforward (see sect.~\ref{sec:misc}).

In the following three sections, the Feynman rules for the correlation
functions of the gauge
field generated by the gradient flow are derived and are
shown to be those of a local field theory with an extra dimension
(the flow time).
The finiteness of the correlation functions
at positive flow time can then be established
using power-counting and the BRS symmetry (sects.~\ref{sec:brs},\ref{sec:fin}),
but for illustration the divergent parts of a set of
one-loop diagrams are first worked out
in sect.~\ref{sec:sample}.
